\subsection{A spherical-equivariantization obstruction for inversion}
\label{subsec:r436-equivariantization}
\begin{theorem}\label{thm:r436-equivariantization}
The fusion ring $\sha{R8_3_436}$ has no complex fusion categorification.
This conclusion uses no near-group associator classification or
unitarization theorem.
\end{theorem}
\begin{proof}
The source tensor has a faithful $C_2$ grading with neutral basis
$[0,1,2,5]$ and odd basis $[3,4,6,7]$. Write $g=X_1$, $X=X_5$ and
$t=X_3$. Its exact products give
\[
g^3=1,\quad gX=Xg=X,\quad X^2=1+g+g^2+3X,
\qquad t^2=1,\quad tgt^{-1}=g^{-1},\quad tXt^{-1}=X.
\]
Suppose $\mathcal C$ categorifies the tensor, and put
$\mathcal D=\mathcal C_e$. The pointed category generated by $t$ may
have a nontrivial associator class $\omega\in H^3(C_2,\C^*)$.
Multiply the associator of the whole graded category by the inflation of
$\omega^{-1}$. This preserves the fusion tensor and the neutral category,
and makes $1,t$ a monoidal section of the grading. Conjugation by that
section gives a coherent categorical $C_2$ action on $\mathcal D$,
interchanging $g,g^2$ and fixing $X$. Coherence is supplied by the section;
it is not inferred from a bare order-two permutation or autoequivalence.

Thornton's sphericality theorem for near-group categories
\cite[Theorem~1.3]{ThorntonNear} applies without a braiding hypothesis.
The spherical dimensions satisfy
\[
\dim g=1,\qquad d=\dim X,\qquad d^2=3+3d.
\]
The pivotal structure on $\mathcal D$ is unique: a tensor automorphism
of the identity has scalars $\lambda_g,\lambda_X$; the relations $gX=X$
and $X\subset X^2$ force both to be one. Every tensor autoequivalence
therefore preserves its spherical structure. The equivariantization
$\mathcal E=\mathcal D^{C_2}$ is consequently spherical, with its
underlying pivotal maps, and the forgetful functor preserves dimensions.

The five simples of $\mathcal E$ are $1,\varepsilon,Y,A,B$, with
forgetful images $1,1,g\oplus g^2,X,X$, respectively. Fixed simples
have two lifts because projective $C_2$ representations over $\C$ can
be linearized. The central sign object $\varepsilon$ fixes $Y$ and
exchanges $A,B$. Forgetful images and Frobenius reciprocity give
\[
Y^2=1+\varepsilon+Y,\qquad YA=AY=YB=BY=A+B.
\]
Since $B=\varepsilon A$, these relations prove commutativity of the
whole ring. If $A^*=B$, the remaining rules would have
$A^2=B^2=\varepsilon+Y+aA+bB$ with $a+b=3$; duality would require
$a=b$, an impossibility. Thus $A,B$ are self-dual and
\[
A^2=B^2=1+Y+aA+bB,\qquad
AB=\varepsilon+Y+bA+aB,\qquad a,b\geq0,\quad a+b=3.
\]
All four ordered pairs are retained. For $s=a-b\in\{-3,-1,1,3\}$,
the ring has scalar character
\[
\chi(1)=1,\quad\chi(\varepsilon)=-1,\quad\chi(Y)=0,
\quad\chi(A)=x,\quad\chi(B)=-x,\qquad x^2=1+sx.
\]
Its formal codegree is
\[
f=2+2x^2=s^2+4\ \pm\ s\sqrt{s^2+4}.
\]
It belongs to $\Q(\sqrt5)\setminus\Q$ for $|s|=1$, and to
$\Q(\sqrt{13})\setminus\Q$ for $|s|=3$. But the spherical dimension
field of $\mathcal E$ is $\Q(\sqrt{21})$, for either root $d$.
Ostrik's formal-codegree field theorem
\cite[Theorem~2.13 and Corollary~2.15]{OstrikPiv3} forbids this: the
corresponding central simple has dimension $\dim(\mathcal E)/f$, whose
forgetful image is an integral combination of the simple dimensions of
$\mathcal E$. Thus $f$ must lie in its dimension field. This contradiction
proves the theorem for arbitrary complex categorifications.
\end{proof}
\begin{remark}[Proof scope and exact certificate]
The previous universal TPE scalar-table parameterization premise is
withdrawn. Its scalar calculation does not prove that missing categorical
statement. The proof above is independent of it and explicitly handles
the pointed associator anomaly. Its source tensor has SHA-256
\par\noindent\sha{0f6d68de67a7287db7dd32a786237f7af9f35ec7c26b3ada2ba50012ff5455c7}.
The checker in \path{evidence/r436_repair_20261024/} verifies the grading,
inversion, every rank-five branch, formal codegrees and the nonintegral
field-norm contradictions. The categorical argument has a separate review.
\end{remark}
